\chapter{Methodology}
\label{chapter:methodology}

\section{Dynamics model and its validation}
\label{sec:dynamics}

The spacecraft state $(\vec{r}, \vec{v}, m)$ evolves under Sun point-mass gravity and a
controllable thrust force $\vec{F}$, with mass loss set by the rocket equation:
\begin{align}
  \dot{\vec r} &= \vec v, &
  \dot{\vec v} &= -\frac{\mu_\odot}{\lVert\vec r\rVert^3}\vec r + \frac{\vec F}{m}, &
  \dot m &= -\frac{\lVert\vec F\rVert}{I_{sp}\,g_0}.
\end{align}
This is integrated with a hand-written fixed-step RK4 scheme (10 substeps per control interval),
using the problem statement's own value of $\mu_\odot$ rather than Tudat's SPICE-derived value, so
that the codebase matches the GTOC4 problem definition it is meant to solve.

To validate this against Tudat's adaptive RKF78 integrator, two checks were run
(\texttt{experiments/validate\_dynamics.py}): one comparing RK4 (problem-statement $\mu_\odot$)
against Tudat (SPICE $\mu_\odot$) over a 300-day coast, and one with both propagators forced to
use the exact same $\mu_\odot$ to isolate the RK4 implementation itself from the choice of gravity
constant. The first comparison showed a residual of \SI{172}{m} in position and \SI{0.043}{mm/s}
in velocity after 300 days; the matched-constant comparison reduced this to \SI{1.9}{m} and
\SI{0.3}{\micro\meter\per\second}. Two candidate explanations for the residual were tested and ruled out or
confirmed empirically rather than assumed:
\begin{itemize}
  \item \textbf{Hypothesis: floating-point/language rounding.} Re-running RK4 at extended
        (80-bit long double) precision shifted the trajectory by only \SI{1.1}{mm} -- two orders
        of magnitude below the observed residual, ruling this out.
  \item \textbf{Confirmed cause: integrator convergence.} Tudat's default relative/absolute
        tolerance of $10^{-10}$ was not fully converged for the eccentric test orbit used;
        bridging a tolerance-$10^{-13}$ Tudat run forward with the (already cross-validated) RK4
        propagator and comparing to the raw tolerance-$10^{-10}$ Tudat output showed a
        $\sim$\SI{640}{m} gap by itself. All downstream Tudat runs in this project therefore use
        tolerance $10^{-13}$.
\end{itemize}
An execution-time comparison (same module) found RK4 to be roughly $8\times$ faster than
Tudat for a representative multi-asteroid propagation -- useful, but far below a
naively-assumed ``$1000\times$'' speedup, since most of Tudat's per-call cost for this problem
size is body/ephemeris setup rather than the integration itself.

\section{Environment}
\label{sec:env}

The task is implemented as a Gymnasium environment (\texttt{src/gtoc4\_env.py}). At each control
step (default: one day) the agent chooses a thrust direction and throttle; the environment
integrates the dynamics forward one control interval and returns a new observation and reward.

\paragraph{Observation (14-D, normalised).} Absolute spacecraft position/velocity
($\vec r/\mathrm{AU}$, $\vec v/v_{ref}$) and state \emph{relative to the target}
($\Delta\vec r/\mathrm{AU}$, $\Delta\vec v/v_{ref}$, where $v_{ref}=\sqrt{\mu_\odot/\mathrm{AU}}$),
plus mass fraction and fraction of time remaining. Both absolute and relative state are included
because gravity depends on absolute position while the rendezvous objective depends on the
relative state.

\paragraph{Action (4-D, $[-1,1]^4$).} A thrust direction in the local radial-transverse-normal
(RTN) frame plus a throttle. RTN separates energy change (transverse), orbit-shape change
(radial) and plane change (normal), which makes the controller analysis in
Section~\ref{sec:wp7} interpretable.

\paragraph{Reward.} Potential-based shaping \parencite{ng1999policy},
$F = \Phi(s') - \Phi(s)$ with $\Phi(s) = -\left(\lVert\Delta\vec r\rVert/\mathrm{AU} +
\lVert\Delta\vec v\rVert/v_{ref}\right)$ and $\gamma=1$ in the shaping term, so that over a full
episode the shaping reward telescopes to net progress made, independent of the path taken. A
terminal bonus scaled by final mass fraction is added on a successful rendezvous.

\paragraph{Termination.} An episode ends on rendezvous (both position and velocity within
tolerance), propellant exhaustion, or divergence ($\lVert\Delta\vec r\rVert > 5\,\mathrm{AU}$).
Running out of time is a \emph{truncation}, not a termination, so PPO's advantage estimator
bootstraps correctly across the time limit rather than treating it as a dead end.

\section{Curriculum}
\label{sec:curriculum}

Three stages of increasing difficulty were used (\texttt{src/curriculum.py}):

\begin{description}
  \item[Stage 1] A synthetic near-Earth target: $\Delta a = 0.02\,\mathrm{AU}$,
    $\Delta i = 0\deg$, $\Delta M = 1\deg$ phase offset, 400-day window. Position/velocity
    tolerance was loosened $10\times$ from an initial \SI{1e5}{km}/\SI{50}{m/s} to
    \SI{1e6}{km}/\SI{500}{m/s} after a diagnostic (Section~\ref{sec:results-stage1}) showed the
    policy reliably closing to within $\sim$\SI{650000}{km}/\SI{200}{m/s} but overshooting.
  \item[Stage 2] A larger synthetic offset: $\Delta a = 0.1\,\mathrm{AU}$, $\Delta i = 2\deg$,
    600-day window, tolerance \SI{1e5}{km}/\SI{100}{m/s}. Never solved despite four different
    fixes (Section~\ref{sec:results-stage2}).
  \item[Stage 3 / randomised] Real GTOC4 asteroids, sampled from a delta-v-feasibility-filtered
    pool (Section~\ref{sec:results-final}), with a 5-day control interval and a randomised
    5--7-year time window.
\end{description}

Both synthetic targets are anchored at the mission start epoch (not propagated forward from the
catalog's far-past reference epoch), to avoid a small $\Delta a$ compounding into tens of degrees
of unintended phase drift by the time the episode begins. The spacecraft starts exactly on
Earth's own heliocentric orbit (no departure kick), so that stage 1's target is genuinely close
to the spacecraft's starting state.

\subsection{Real-asteroid reachability}
\label{sec:reachability}

The original plan called for a 300--600-day window for the real-asteroid stages. Scanning the
full 1436-asteroid catalog for the closed-form impulsive delta-v needed to reach each asteroid
(an energy term from the semi-major-axis difference, plus a phasing term $v\,\Delta\theta$ from
the initial velocity-direction mismatch) showed that only 2 asteroids are reachable at all within
600 days, and none with a comfortable margin, against a 600-day thrust budget of
$a_{max}\cdot\Delta t \approx$ \SI{4666}{m/s}. This is not a training artifact: it explains why
GTOC4's own mission budget is 10 years, not 1.5. Extending the window to 5--7 years yields 17--37
reachable asteroids with a $1.5\times$ delta-v margin, which is the pool used for stage 3 and the
randomised-target configuration.

\section{Reinforcement learning setup}
\label{sec:rl-setup}

Training uses Stable-Baselines3's \parencite{raffin2021sb3} PPO implementation, with an MLP policy (two hidden layers of 64 units, tanh
activation), learning rate $3\times10^{-4}$, $n_{steps}=2048$, batch size 64, 10 epochs per
update, $\gamma=0.99$, GAE $\lambda=0.95$, clip range 0.2, no entropy bonus. Observations are
normalised online (\texttt{VecNormalize}, observation-only, reward left unnormalised while
shaping was being tuned). Two baselines -- coasting (zero thrust) and constant tangential thrust
-- are used throughout as reference points; a policy that cannot beat coasting has not learned
anything about steering towards the target. Where noted, later curriculum stages are
\emph{warm-started} from an earlier stage's trained weights.
